% Appendices

\chapter*{Appendices}
\addcontentsline{toc}{chapter}{Appendices}

\section*{Appendix A: System Architecture Diagram}

[Insert system architecture diagram here]

\section*{Appendix B: Database Schema}

[Insert complete database schema diagram here]

\section*{Appendix C: API Documentation}

\subsection*{API Endpoints}

\subsubsection*{POST /api/register\_farmer}

Registers a new farmer in the system.

\textbf{Request Body:}
\begin{verbatim}
{
  "full_name": "John Doe",
  "district": "Serowe"
}
\end{verbatim}

\textbf{Response:}
\begin{verbatim}
{
  "status": "success",
  "owner_id": "OWN-ABC12345"
}
\end{verbatim}

\subsubsection*{POST /api/ingest\_telemetry}

Ingests telemetry data from sensors.

\textbf{Request Body:}
\begin{verbatim}
{
  "cattle_id": "BOT-001",
  "temperature": 38.5,
  "heart_rate": 72
}
\end{verbatim}

\textbf{Response:}
\begin{verbatim}
{
  "status": "success",
  "validation": "valid",
  "orchestrator_result": "[Thekiso Agent] READY: BOT-001 qualifies for market liquidation"
}
\end{verbatim}

\section*{Appendix D: Test Cases}

\subsection*{Unit Test Cases}

\subsubsection*{Temperature Validation}

\begin{table}[H]
\centering
\begin{tabular}{|l|l|l|l|}
\hline
\textbf{Test ID} & \textbf{Input} & \textbf{Expected} & \textbf{Actual} \\
\hline
TC-001 & 38.5°C & Valid & Valid \\
TC-002 & 28.0°C & Invalid & Invalid \\
TC-003 & 47.5°C & Invalid & Invalid \\
TC-004 & 30.0°C & Valid & Valid \\
TC-005 & 45.0°C & Valid & Valid \\
\hline
\end{tabular}
\caption{Temperature Validation Test Cases}
\end{table}

\section*{Appendix E: Code Samples}

\subsection*{Validation Engine Code}

\begin{lstlisting}[language=Python, caption=Temperature Validation]
def validate_temperature(self, temp):
    """Validate temperature is within physiological bounds"""
    # Added extra check because sensor was sending -0.1 sometimes
    if temp < self.MIN_TEMP or temp > self.MAX_TEMP:
        return False, f"Temperature {temp}C outside valid range ({self.MIN_TEMP}-{self.MAX_TEMP}C)"
    return True, "Valid"
\end{lstlisting}

\subsection*{Graph Orchestrator Code}

\begin{lstlisting}[language=Python, caption=Self-Healing Implementation]
def route_telemetry(self, cattle_id):
    try:
        # Fetch verified data from our INFS-validated database
        conn = sqlite3.connect(self.db_path)
        cursor = conn.cursor()
        cursor.execute("SELECT temperature FROM telemetry_logs 
                        WHERE cattle_id = ? ORDER BY log_id DESC LIMIT 1", 
                        (cattle_id,))
        row = cursor.fetchone()
        conn.close()

        if not row:
            # COMP: Self-healing fallback to neutral supervisor node
            return self.agent_supervisor(cattle_id, "NO_DATA_FOUND")
        
        # ... rest of implementation
    except Exception as e:
        # COMP: Self-healing - catch exception and route to supervisor
        print(f"[COMP] Exception caught: {e}, routing to supervisor")
        return self.agent_supervisor(cattle_id, f"PROCESSING_ERROR: {str(e)}")
\end{lstlisting}

\section*{Appendix F: User Manual}

\subsection*{Getting Started}

1. Install Python 3.10 or higher
2. Install required dependencies: \texttt{pip install -r requirements.txt}
3. Initialize database: \texttt{python init\_db.py}
4. Start the application: \texttt{python app.py}
5. Access the web interface at \texttt{http://localhost:5000}

\subsection*{Registering as a Farmer}

1. Navigate to the registration page
2. Enter your full name and district
3. Submit the registration form
4. Note your generated owner ID for future reference

\subsection*{Registering Cattle}

1. Log in with your credentials
2. Navigate to the Cattle Management section
3. Click "Register New Cattle"
4. Enter cattle details (ID, breed, birth date, gender)
5. Submit the form

\subsection*{Viewing Telemetry Data}

1. Navigate to the Dashboard
2. Select a cattle from your list
3. View current health status and telemetry history
4. Review agent recommendations

\section*{Appendix G: Installation Guide}

\subsection*{System Requirements}

\begin{itemize}
    \item Operating System: Windows 10/11, Linux, or macOS
    \item Python: 3.10 or higher
    \item RAM: 4GB minimum (8GB recommended)
    \item Disk Space: 500MB minimum
    \item Internet: Required for initial setup and API calls
\end{itemize}

\subsection*{Installation Steps}

\begin{enumerate}
    \item Clone the repository: \texttt{git clone https://github.com/PipsWhisperer/LesakaAi\_project.git}
    \item Navigate to project directory: \texttt{cd LesakaAi\_project}
    \item Create virtual environment: \texttt{python -m venv venv}
    \item Activate virtual environment:
    \begin{itemize}
        \item Windows: \texttt{venv\textbackslash Scripts\textbackslash activate}
        \item Linux/macOS: \texttt{source venv/bin/activate}
    \end{itemize}
    \item Install dependencies: \texttt{pip install -r requirements.txt}
    \item Initialize database: \texttt{python init\_db.py}
    \item Start application: \texttt{python app.py}
\end{enumerate}

\section*{Appendix H: Glossary}

\begin{description}
    \item[3NF] Third Normal Form - A database design principle that eliminates redundancy
    \item[API] Application Programming Interface - Set of protocols for building software
    \item[DCG] Directed Cyclic Graph - A graph with directed edges that may contain cycles
    \item[INFS] Information Systems - Academic module focusing on data governance
    \item[IoT] Internet of Things - Network of interconnected physical devices
    \item[RBAC] Role-Based Access Control - Method of regulating access to computer resources
    \item[COMP] Computer Science - Academic module focusing on algorithms and systems
    \item[REST] Representational State Transfer - Architectural style for web services
    \item[SQLite] Relational database management system contained in a C library
\end{description}
